% Modernes Code-Highlighting mit minted für LuaLaTeX
% Ersetzt das veraltete listings-Paket komplett

% --- Hyperref-kompatible Minted-Konfiguration ---
% minted lädt hyperref automatisch - wir konfigurieren es nach dem Laden
\PassOptionsToPackage{
  unicode=true,
  breaklinks=true,
  colorlinks=true
}{hyperref}

% --- Minted-Paket mit LuaLaTeX-Optimierungen ---
\usepackage[
  cache=true,                    % Caching für bessere Performance
  cachedir=build/minted-cache,   % Cache-Verzeichnis
  frozencache=false              % Cache aktualisierbar
]{minted}

% --- Moderne Farbdefinitionen (konsistent mit Template) ---
% Nutzt die bereits definierten Farben aus Farben_definieren-settings.tex
\definecolor{codebg}{RGB}{248, 248, 248}        % Hintergrund
\definecolor{codeborder}{RGB}{224, 224, 224}    % Rahmen

% --- Standard-Minted-Konfiguration ---
\setminted{
  % === Layout ===
  bgcolor=codebg,                % Hintergrundfarbe
  frame=single,                  % Einfacher Rahmen
  framesep=10pt,                 % Abstand zum Rahmen
  rulecolor=codeborder,          % Rahmenfarbe
  % === Schrift ===
  fontsize=\small,               % Schriftgröße
  fontfamily=tt,                 % Monospace
  % === Verhalten ===
  breaklines=true,               % Automatischer Zeilenumbruch
  breakanywhere=true,            % Umbruch überall möglich
  autogobble=true,               % Automatisches Einrückung entfernen
  % === Nummerierung ===
  linenos=false,                 % Standardmäßig keine Zeilennummern
  xleftmargin=0pt,               % Kein extra Rand
  % === Moderne Features ===
  style=default,                 % Standard-Highlighting-Style
  tabsize=2                      % Tab-Größe
}

% --- Sprachspezifische Konfigurationen ---
% Für besonders häufig verwendete Sprachen

\setminted[python]{
  style=default,
  python3=true,
  linenos=true,
  numbersep=8pt
}

\setminted[scala]{
  style=default,
  linenos=true,
  numbersep=8pt
}

\setminted[java]{
  style=default,
  linenos=true,
  numbersep=8pt
}

\setminted[javascript]{
  style=default,
  linenos=true,
  numbersep=8pt
}

\setminted[rust]{
  style=default,
  linenos=true,
  numbersep=8pt
}

\setminted[json]{
  style=default,
  linenos=false
}

\setminted[yaml]{
  style=default,
  linenos=false
}

% --- Moderne Themes (auskommentiert - aktivierbar) ---
% Verschiedene Highlighting-Styles:
% \setminted{style=github}        % GitHub-Style
% \setminted{style=vs}            % Visual Studio
% \setminted{style=monokai}       % Monokai (dunkel)
% \setminted{style=solarized}     % Solarized
% \setminted{style=colorful}      % Farbig

% --- Convenience-Befehle ---
% Inline-Code (ersetzt \lstinline)
\newcommand{\inlinecode}[2][text]{\mintinline{#1}{#2}}

% Code aus Datei einbinden (ersetzt \lstinputlisting)
\newcommand{\inputcode}[3][]{%
  \inputminted[#1]{#2}{#3}%
}

% Spezifische Sprach-Befehle für häufige Verwendung
\newcommand{\pythoninput}[2][]{\inputminted[linenos,#1]{python}{#2}}
\newcommand{\javainput}[2][]{\inputminted[linenos,#1]{java}{#2}}
\newcommand{\jsinput}[2][]{\inputminted[linenos,#1]{javascript}{#2}}
\newcommand{\rustinput}[2][]{\inputminted[linenos,#1]{rust}{#2}}

% --- Moderne Code-Umgebungen ---

% Standard-Code-Umgebung
\newenvironment{moderncode}[2][]{%
  \VerbatimEnvironment
  \begin{minted}[#1]{#2}%
}{%
  \end{minted}%
}

% Code mit Zeilennummern
\newenvironment{numberedcode}[2][]{%
  \VerbatimEnvironment
  \begin{minted}[linenos,#1]{#2}%
}{%
  \end{minted}%
}

% --- Caption-Integration ---
% Für Code-Listings mit Beschriftungen
\usepackage{newfloat}
\DeclareFloatingEnvironment[
  fileext=lol,
  listname={Codeverzeichnis},
  name=Code,
  placement=tbp,
]{codefloat}

% Code-Float-Umgebung
\newenvironment{codelisting}[3][htbp]{%
  \begin{codefloat}[#1]
  \caption{#3}
  \begin{minted}{#2}%
}{%
  \end{minted}
  \end{codefloat}%
}

% --- Unterstützte Sprachen (300+) ---
% Beispiele der wichtigsten:
% - Python, Java, JavaScript, TypeScript
% - Rust, Go, Swift, Kotlin, Scala, Dart
% - C, C++, C#, PHP, Ruby, Perl
% - HTML, CSS, XML, JSON, YAML, TOML
% - SQL, bash, PowerShell, Dockerfile
% - LaTeX, Markdown, und viele mehr!

% --- Performance-Optimierungen ---
% Shell-Escape aktivieren (automatisch von minted)
\makeatletter
\@ifpackageloaded{shellesc}{}{%
  \PackageWarning{minted-settings}{%
    Shell-Escape nicht aktiviert. Verwenden Sie --shell-escape beim Kompilieren.%
  }
}
\makeatother

% --- Debug-Information ---
\AtBeginDocument{%
  \typeout{=== Minted Code-Highlighting ===}
  \typeout{Engine: Pygments (300+ Sprachen)}
  \typeout{Cache: aktiviert (build/minted-cache)}
  \typeout{Features: Automatischer Umbruch, Unicode}
  \typeout{Kompilierung: lualatex --shell-escape}
  \typeout{Beispiel: \string\begin{minted}{python}...\string\end{minted}}
  \typeout{===================================}
}

% --- Verwendungsbeispiele (auskommentiert) ---
% \begin{minted}{python}
% def fibonacci(n):
%     if n <= 1:
%         return n
%     return fibonacci(n-1) + fibonacci(n-2)
% \end{minted}
%
% \begin{minted}{rust}
% fn main() {
%     println!("Hello, world! ��");
% }
% \end{minted}
%
% \inputcode{python}{scripts/example.py}
% \inlinecode[python]{print("inline code")}
%
% \begin{codelisting}{javascript}{Moderne JavaScript-Funktion}
% const api = async () => {
%   const response = await fetch('/api/data');
%   return response.json();
% };
% \end{codelisting}